\chapter{Desarrollo de la propuesta y trabajo realizado}

\section{Preparacion del entorno}
Esta sección recogerá la configuración real del entorno: sistema operativo, estructura de repositorios, despliegue mediante contenedores, red interna, almacenamiento, GPU, drivers y herramientas de documentación. Debe incluir comandos relevantes, capturas y decisiones tomadas durante la puesta en marcha.

\section{Fase 1: infraestructura base}
\placeholder{Describir instalación del servidor, particionado/almacenamiento, Docker o motor equivalente, redes, volúmenes persistentes, políticas de reinicio y estructura de backups.}

\section{Fase 2: inferencia local}
\placeholder{Profundizar en Ollama, llama.cpp u otras herramientas utilizadas: instalación, modelos probados, cuantización, uso de VRAM, latencia, problemas con AMD/ROCm, fallback y resultados.}

\section{Fase 3: orquestación y automatización}
\placeholder{Explicar n8n, Node-RED u orquestador final: flujos creados, nodos, entradas, salidas, gestión de errores, integración con IA y servicios internos.}

\section{Fase 4: domótica e interacción}
\placeholder{Describir Home Assistant, dispositivos integrados, satélites de voz, automatizaciones domésticas y casos de interacción natural.}

\section{Fase 5: nube privada y memoria}
\placeholder{Explicar servicios de almacenamiento, notas, documentos, calendarios, RAG, embeddings, base vectorial, indexación y recuperación de conocimiento.}

\section{Fase 6: seguridad, acceso remoto y monitorización}
\placeholder{Detallar VPN, proxy inverso si aplica, certificados, segmentación, alertas, dashboards, métricas, healthchecks, limpieza de caché y autoreparación.}

\section{Problemas encontrados}
\placeholder{Documentar incidencias reales: drivers, memoria, latencia, conflictos entre contenedores, permisos, seguridad, backups, integraciones inestables o recortes de alcance.}
